\chapter{Desarrollo y Consignas}

\section{Implementación del Amplificador}
Implementar un amplificador realimentado según la topología del circuito de la Figura \ref{fig:esquematico}.

\vspace{1cm}
\begin{figure}[H]
    \centering
    % \includegraphics[width=0.8\textwidth]{images/esquematico.png}
    % \caption{Circuito esquemático del Amplificador de 2 etapas.}
    % \label{fig:esquematico}

    \begin{tikzpicture}[transform shape]
        % Paths, nodes and wires:
        \draw (2, 4) to[american resistor, l={$R_{B3}$}] (2, -0);
        \draw (2, -0) to[american resistor, l={$R_{B4}$}] (2, -5);
        \node[npn](N1) at (4, -0){} node[anchor=west] at (N1.text){$Q_2$};
        \draw (4, 4) to[american resistor, l={$R_{C2}$}] (4, 2);
        \draw (4, -0.77) to[american resistor, l={$R_{E3}$}] (4, -3);
        \draw (4, -3) to[american resistor, l={$R_{E4}$}] (4, -5);
        \draw (4, -3) -| (5.5, -3.5);
        \draw (4, 2) -| (4, 0.77);
        \draw (5.5, -5) -- (2, -5);
        \node[npn](N2) at (-1, -0){} node[anchor=west] at (N2.text){$Q_1$};
        \draw (-1, 4) to[american resistor, l={$R_{C1}$}] (-1, 2);
        \draw (-1, 2) -| (-1, 0.77);
        \draw (1, -0) -| (0.5, 1.5) -- (-1, 1.5);
        \draw (-1, -0.77) to[american resistor, l={$R_{E1}$}] (-1, -3);
        \draw (-1, -3) to[american resistor, l={$R_{E2}$}, label distance=-0.06cm] (-1, -5);
        \draw (2, -5) -- (-1, -5);
        \draw (-1, -3) -| (0.4, -3.5);
        \draw (-3, 4) to[american resistor, l={$R_{B1}$}] (-3, -0);
        \draw (-3, -0) to[american resistor, l={$R_{B2}$}] (-3, -5);
        \draw (-3, -0) -- (-1.84, -0);
        \draw (-3, 4) -- (4, 4);
        \draw (-3, -5) -- (-1, -5);
        \draw (0.4, -3.5) to[curved capacitor, l={$C_{E1}$}, label distance=-0.07cm] (0.4, -5);
        \draw (5.5, -3.5) to[curved capacitor, l={$C_{E2}$}, label distance=0.01cm] (5.5, -5);
        \draw (-3, -0) to[curved capacitor, l={$C_{a1}$}] (-4.75, 0);
        \draw (-6, -2) to[sinusoidal voltage source, l={$V_S$}] (-6, -3);
        \draw (-6, -2) -| (-6, -0) -- (-4.75, -0);
        \draw (-6, -3) -| (-6, -5) -- (-3, -5);
        \draw (-7.5, 1) to[battery, l={$V_{CC}$}] (-7.5, -1);
        \draw (-7.5, 1) -| (-7.5, 4) -- (-3, 4);
        \draw (-7.5, -1) |- (-6, -5);
        \draw (4, 1.5) -- (6, 1.5);
        \draw (6.75, 1.5) to[curved capacitor, l={$C_0$}] (8, 1.5);
        \draw (8, 1.5) to[american resistor, l={$R_L$}] (8, -5);
        \draw (7.5, -5) -- (5.5, -5);
        \draw (5.6, -6.2) to[curved capacitor, l_={$C_f$}] (4.1, -6.2);
        \draw (6, 1.5) -- (6.25, 1.5);
        \draw (6.75, 1.5) |- (6.3, -6.2);
        \draw (1.6, -6.2) |- (-1, -0.77);
        \node[circ] at (-3, -0){};
        \node[circ] at (-1, 1.5){};
        \node[circ] at (-1, 4){};
        \node[circ] at (-3, 4){};
        \node[circ] at (-3, -5){};
        \node[circ] at (-6, -5){};
        \node[circ] at (-1, -5){};
        \node[circ] at (-1, -3){};
        \node[circ] at (0.4, -5){};
        \node[circ] at (2, -5){};
        \node[circ] at (4, -5){};
        \node[circ] at (4, -3){};
        \node[circ] at (5.5, -5){};
        \node[circ] at (6.75, 1.5){};
        \node[circ] at (4, 1.5){};
        \node[circ] at (2, 4){};
        \node[circ] at (-1, -0.77){};
        \draw (8, -5) -- (7.5, -5);
        \draw (6.3, -6.2) -- (5.8, -6.2);
        \draw (6.25, 1.5) -- (6.75, 1.5);
        \node[ground] at (7.5, -5){};
        \node[circ] at (7.5, -5){};
        \draw (1, -0) to[curved capacitor, l={$C_{a2}$}, label distance=-0.07cm] (2, -0);
        \draw (5.8, -6.2) -- (5.6, -6.2);
        \draw (4.1, -6.2) to[american resistor, l_={$R_f$}] (1.6, -6.2);
    \end{tikzpicture}
    \caption{Circuito esquemático del Amplificador de 2 etapas.}
    \label{fig:esquematico}
\end{figure}

\subsection{Valores de Componentes Normalizados}
A continuación se detallan los valores de los componentes a utilizar para el armado del circuito:

\begin{table}[!ht]
    \centering
    \begin{tabular}{|l|l||l|l||l|l||l|l|}
        \hline
        $V_{cc}$ & $22 \V$    & $R_{e1}$ & $130 \ohm$  & $R_{c2}$ & $2.0 \kohm$ & $R_f$    & $1.6 \kohm$ \\ \hline
        $R_{b1}$ & $820 \kohm$ & $R_{e2}$ & $12 \kohm$   & $R_{e3}$ & $240 \ohm$  & $C_{a1}$ & $1 \uF$    \\ \hline
        $R_{b2}$ & $560 \kohm$ & $R_{b3}$ & $110 \kohm$  & $R_{e4}$ & $200 \ohm$  & $C_{E1}$ & $10 \uF$   \\ \hline
        $R_{c1}$ & $18 \kohm$  & $R_{b4}$ & $24 \kohm$   & $R_L$    & $470 \ohm$  & $C_{a2}$ & $1 \uF$    \\ \hline
        $C_{e2}$ & $10 \uF$   & $C_o$    & $1 \uF$     & $Q_1$    & BC 337      & $Q_2$    & BC 337     \\ \hline
    \end{tabular}
    \caption{Valores de Componentes Normalizados}
    \label{tab:componentes}
\end{table}

\section{Análisis Teórico y Cálculos}
[COMPLETAR CON LOS CÁLCULOS TEÓRICOS DEL PUNTO DE OPERACIÓN Y PARÁMETROS DINÁMICOS]
